\documentclass[12pt]{article}
\usepackage[T1]{fontenc}
\usepackage[utf8]{inputenc}
\usepackage{lmodern}
\usepackage{textcomp}
\usepackage{enumitem}
\usepackage{geometry}
\geometry{margin=1in}
\begin{document}
\begin{center}
Answer Key - Astronomy - Graduate Level Assessment 11 \
\small (Answers are ordered exactly as in the question paper.)
\end{center}

\section*{Multiple Choice}
\begin{enumerate}[label=\arabic*.]
\item \textbf{Answer:} D
\\ \textit{Options:} A) Complex organic molecules specifically PAHs; B) Magnetite grains similar to those formed by bacteria on Earth; C) Carbonate minerals indicating a thicker warmer Martian atmosphere; D) Amino acids with a preferred orientation or "chirality"
\\ \textit{Note:} Amino acids with a preferred orientation or "chirality" were not cited as evidence for life in the martian meteorite ALH 84001 because while certain organic compounds were found, the specific chiral preference of amino acids, which is indicative of biological processes, was not conclusively demonstrated in the meteorite.
\item \textbf{Answer:} D
\\ \textit{Options:} A) ExoMars; B) Mars Exploration Rovers; C) Mars Science Laboratory; D) Phoenix Mars Lander
\\ \textit{Note:} The Phoenix Mars Lander is specifically designed to explore the Martian arctic region and investigate the presence of water ice by digging into the soil, making it the mission that aligns with the details provided in the question.
\item \textbf{Answer:} B
\\ \textit{Options:} A) You see only the star's blackbody spectrum.; B) You see the star's blackbody spectrum with absorption lines due to hydrogen.; C) You see only emission lines characteristic of hydrogen.; D) You see only emission lines characteristic of the star's composition.
\\ \textit{Note:} The correct answer is accurate because the cooler hydrogen gas absorbs specific wavelengths of light from the star's blackbody spectrum, resulting in dark absorption lines at those wavelengths in the observed spectrum.
\item \textbf{Answer:} A
\\ \textit{Options:} A) The U.S.S.R.; B) The U.S.; C) France; D) A and B
\\ \textit{Note:} The U.S.S.R. successfully deployed several landers to Venus, including the Venera series, which provided invaluable data about the planet's atmosphere and surface conditions, making it the only country to have achieved this feat.
\item \textbf{Answer:} C
\\ \textit{Options:} A) The azimuth of the NCP is the angular distance from the North Pole.; B) The azimuth of the NCP is the same as your latitude.; C) The altitude of the NCP is the same as your latitude.; D) The altitude of the NCP is your angular distance from the North Pole.
\\ \textit{Note:} The altitude of the north celestial pole directly corresponds to an observer's latitude because as one moves north or south on Earth, the angle of the NCP above the horizon changes correspondingly, making it a reliable indicator of one's geographic position.
\end{enumerate}

\section*{True / False}
\begin{enumerate}[label=\arabic*.]
\setcounter{enumi}{5}
\item \textbf{Answer:} False
\\ \textit{Options:} A) True; B) False
\\ \textit{Note:} The Large Magellanic Cloud is a dwarf galaxy located approximately 163,000 light-years away from Earth, whereas the closest planetary nebula, NGC 3918, is located only about 4,500 light-years away.
\item \textbf{Answer:} True
\\ \textit{Options:} A) True; B) False
\\ \textit{Note:} Type Ia supernovae are indeed the result of a white dwarf in a binary system accumulating material from a companion star, leading to a thermonuclear explosion when it reaches the Chandrasekhar limit of approximately 1.4 solar masses.
\item \textbf{Answer:} False
\\ \textit{Options:} A) True; B) False
\\ \textit{Note:} The length of a year is determined by the Earth's orbit around the Sun, which is influenced by the gravitational force; if the gravitational force were to change, the Earth's orbital characteristics, including the length of a year, would also change.
\item \textbf{Answer:} False
\\ \textit{Options:} A) True; B) False
\\ \textit{Note:} Sunspots actually appear in cycles approximately every 11 years, known as the solar cycle, rather than every 9 years.
\item \textbf{Answer:} False
\\ \textit{Options:} A) True; B) False
\\ \textit{Note:} The statement is false because the correct order of electromagnetic radiation from shortest to longest wavelength is gamma rays, X rays, ultraviolet, visible light, infrared, and radio waves.
\end{enumerate}

\section*{Long Form Response}
\begin{enumerate}[label=\arabic*.]
\setcounter{enumi}{10}
\item \textbf{Answer:} Gravitational resonance with Jupiter plays a critical role in inhibiting the formation of a planet in the asteroid belt by creating a dynamic environment that disrupts the accumulation of material. The gravitational influence of Jupiter causes periodic perturbations in the orbits of asteroids, which prevents them from coalescing into a larger body. This phenomenon, often referred to as the "Jupiter barrier," results in a stable population of small bodies rather than a singular planet. The implications of this resonance extend to our understanding of planetary formation, as it underscores the significance of gravitational interactions among celestial bodies in shaping planetary systems and highlights the complexities involved in the processes of accretion and stability in the early solar system. Thus, the relationship between Jupiter and the asteroid belt illustrates how external gravitational forces can dictate the evolutionary pathways of forming planets.
\\ \textit{Note:} correct because it accurately describes how Jupiter's gravitational resonance disrupts the orbits of asteroids in the belt, preventing them from merging into a planet, and highlights the broader implications for understanding the conditions necessary for planetary formation.
\item \textbf{Answer:} The current tilt of the Mars Exploration Rover Spirit towards the north can be primarily attributed to its location in the southern hemisphere during the winter season. In this hemisphere, solar insolation diminishes significantly, leading to a decrease in energy absorption by the rover's solar panels. As the rover attempts to maximize its exposure to the sun for energy, its orientation naturally shifts northward. Additionally, the seasonal changes on Mars can cause variations in surface conditions, influencing the rover's stability and tilt. Thus, the tilt towards the north is a strategic adaptation to the environmental challenges posed by the Martian winter, allowing Spirit to sustain its operations despite the harsh conditions.
\\ \textit{Note:} The answer correctly identifies that the Mars Exploration Rover Spirit tilts northward during the southern hemisphere winter to optimize solar panel exposure due to reduced solar insolation, reflecting the rover's adaptive response to seasonal energy challenges.
\end{enumerate}

\vfill
\noindent\textit{End of Answer Key}
\end{document}